\chapter*{Anexo}
\addcontentsline{toc}{chapter}{Anexo}
\label{cap:capitulo8}

\section*{Anexo I: Repositorio del proyecto}
\addcontentsline{toc}{section}{Anexo I: Repositorio del proyecto}

Con el objetivo de facilitar la reproducción del trabajo desarrollado, además de poder consultar todo el código fuente implementado durante este Trabajo de Fin de Grado, éste se encuentra disponible en un repositorio de GitHub\footnote{\url{https://docs.github.com/es}} de forma totalmente pública.

\vspace{0.25cm}

Este repositorio incluye la implementación completa del modelo cinemático desarrollado, la arquitectura software basada en ROS 2\footnote{\url{https://docs.ros.org/en/jazzy/index.html}} en la que se ha integrado, los paquetes utilizados para la generación de trayectorias, el modelo URDF\footnote{\url{https://docs.ros.org/en/jazzy/Tutorials/Intermediate/URDF/URDF-Main.html}} empleado para la simulación, ficheros de configuración como el de RViZ\footnote{\url{https://docs.ros.org/en/jazzy/Tutorials/Intermediate/RViz/RViz-Main.html}}, además de otros ejemplos de ejecución de pruebas de validación iniciales que completan la estructura del proyecto.

\vspace{0.25cm}

El hecho de que este repositorio sea público para todo el mundo, permite que otros desarrolladores puedan reutilizarlo para ayudar a ampliar el sistema desarrollado, facilitando así tanto la validación de los resultados obtenidos como la evolución del proyecto hacia una futura implementación física del robot.

\vspace{0.25cm}

Este repositorio puede consultarse accediendo al siguiente enlace\footnote{\url{https://github.com/aleon2020/cable\_driven\_parallel\_robot\_2d}}.
